\begin{proof}[Proof of \expref{Theorem}{prop2}].   %
Here, intending to recycle $n$, $m$ and $f$, we rename
these quantities
$\nn$, $\mm$ and $\ff$.
We may of course assume $\gd$ is fairly small.
Let (for example) $\xi = \gd/3$, fix $\eps$ with
$0<\eps < \xi^2$, and set
$k=(1+\gd)\log \nn$ and $\mm = \eps 2^k \nn$.
These values are easily seen to give the hypotheses of
\expref{Lemmas}{mSlemma} \expref{and}{muflemma}.
In particular, we can say that w.\,h.\,p.
the supports of the %
$C_i$  %
are chosen so \eqref{mS} holds
(note this says nothing about the values specified by
%
the $g_i$) %
and
\beq{mS'}
\mathbb E (\ff)\sim (1-2^{-k})^\mm\sim \eee^{-\eps \nn}.
\enq




Set $n=(1/2+\gd) \nn$.
Fix %
$S\in \binom{[\nn]}{s}$, %
set $m=m_S$,
and
let $f=f_S$ be the %
conjunction  %
of the $m$ %
clauses $C_i$%
---w.l.o.g. %
$C_1,\dots, C_m$%
---that miss $S$.
Thus $f$ is the %
conjunction of %
$m\sim (1/2+\gd)^k\mm =\eps (1+2\gd)^k\nn$
random $k$-clauses from a universe of $n$ variables.
\expref{Theorem}{prop2} (with $\ga=\gd$) thus follows from

\mn
\emph{Claim A.}
$\Pr(\mu(f)  > \exp[-n^{\gd}]) < o(2^{-\nn})$

\mn
(since then w.\,h.\,p. we have $\mathbb E (f_S)\leq \exp[-n^{\gd}]$ for every $S$).


\mn
\emph{Remarks.}
The actual bound in Claim A will be
$\exp[-\gO(m)]$,
so much smaller than $2^{-\nn}$.
Note that here it %
does not %
matter whether we take $\mu$ to be
our original measure (\ie, $\mu$ uniform on $\{0,1\}^\nn$)
or uniform measure on $Q:=\{0,1\}^n$;
but it is %
now more natural to think of
the latter---and we will do so in what follows---since
our original universe plays no further role in this discussion.
It may also be worth noting that, unlike in the proof of
\expref{Lemma}{muflemma}, the second moment method is not strong
enough to give the exponential bound in Claim A.




\mn
\emph{Claim B.}
If $X\sub Q$,
$\mu (X) =\gb> \exp[-o(n/\log^2n)]$ and $\gz =o(2^{-k})$,
then for a random $k$-clause $C$,
$$
\Pr(\mu (C\wedge X) >(1-\gz)\mu (X)) < 1/2
$$
(where $C\wedge X =\{x\in X: C\sim x\}$).


\mn
\emph{Remark.}  This is probably true for
$\gb$ greater than something like $\exp[-n/k]$.
The bound in the claim
is just what the proof gives, and is more than enough
for us since %
we are %
really interested in much larger $\gb$.


\medskip
To see that Claim B implies Claim A,
set $f_j=\bigwedge_{i=1}^jC_i$ and notice that $\mu(f) \geq \gb$
implies (for example)
\beq{imufi}
%
%
\left|\left\{i: \mathbb E (f_i) <\left(1-\tfrac{5\ln(1/\gb)}{m}\right)
\mathbb E (f_{i-1})\right\}\right| < m/5
%
\enq
(and, of course, $\mathbb E (f_i)\geq \gb$ for all $i$).
But if we take
$\gb =\exp[-n^{\gd}]$ then our choice of parameters gives
$$\gz:=5m^{-1}\ln(1/\gb) =o(2^{-k})$$
(using $m/\ln (1/\gb) =
\Theta(n^{(1+\gd)\log (1+2\gd)+1-\gd})$ and
$2^k = n^{1+\gd}$),
so Claim B bounds the probability of \eqref{imufi} by
%
\begin{equation*}
\Cc{m}{m/5} 2^{-4m/5} =o(2^{-\nn}).   \qedhere    %
\end{equation*}
%
\end{proof}   %

\mn
\emph{Proof of Claim B.}
Let $G$ be the bipartite graph on
$Q\cup W$, where $W$ is the set of $(n-k)$-dimensional
subcubes of $Q$ and, for $(x,D)\in Q\times W$,
we take $x\sim D$ if $x\in D$.
(So %
we have  %
gone to complements:  for a clause $C$
the corresponding subcube is $D=\{y:C\not\sim y\}$,
so $C\wedge X =X\sm D$.)

Assuming Claim B fails at $X$, fix $\SSS \sub W$ with $|\SSS |=|W|/2$ and
$$D\in \SSS
  ~\Ra ~\mu (D\cap X) < \gz\mu(X),$$
and set $\TTT =W\sm \SSS $.

%
Consider the experiment:  %
(i) choose $x$ uniformly from $X$;
(ii)  choose $D$ uniformly from the members of $W$ containing $x$;
(iii)  choose $y$ uniformly from $D$.


\mn
\emph{Claim C.}  $\Pr(y\in X)> (2-o(1))\gb$.


\mn
\begin{proof}   %
%
Since each triple $(x,D,y)$ with $x\in X$ and $x,y\in D$ is
produced by (i)-(iii) with probability
$|X|^{-1}\cdot 2^k|W|^{-1}\cdot 2^{k-n}$,
we just need to show that the number of such triples with
$y\in X$ is at least
$$
(2-o(1))\gb |X||W|2^n2^{-2k}  = (2-o(1))|X|^2|W|2^{-2k}.
$$
Writing $d$ for degree in $G$, we have
$$
\sum_{x\in X}d_\SSS(x) =\sum_{D\in \SSS }d_X(D) < |\SSS |\gz
 |X|,
$$
implying
%
\begin{align*}
\sum_{D\in \TTT }d_X(D)& = 
\sum_{x\in X}(d(x)-d_\SSS (x))\\
&> |X||W|2^{-k} - \gz |\SSS ||X| = (1-o(1))|X||W|2^{-k}.
\end{align*}   %
The number of %
triples $(x,D,y)$  %
as above is thus
%
\begin{align*}
\sum_{D\in W} d_X^2(D)
&\geq
\sum_{D\in \TTT } d_X^2(D)
\geq
\left(\sum_{D\in \TTT } d_X(D)\right)^2/|\TTT |\\   %
&> (1-o(1))|X|^2|W|^2|\TTT |^{-1}2^{-2k}
= (2-o(1))|X|^2|W|2^{-2k}\,.
 \qedhere                          
\end{align*}   %
%
\end{proof}   %

Let $T(x)$ be
the random element of $Q$ %
obtained %
from $x$
by choosing $K$ uniformly from $\binom{[n]}{k}$ and
randomly (uniformly, independently) 
resampling %
the %
$x_i$  %
%
%
%
for $i\not\in K$.    %
Then $y$ %
obtained  %
from $x$ by (ii) and (iii) above is just $T(x)$,
so
the next assertion contradicts Claim C, completing the proof of
Claim B (and \expref{Theorem}{prop2}).



\mn
\emph{Claim D.}
If $\mu(X) > \exp[-o(n/\log^2n)]$
and $x$ is uniform from $X$, then $\Pr(T(x)\in X) < (1+o(1))\mu(X)$.

\mn
\emph{Remark.}
If $X$ is a subcube of codimension $n/k$, say
$X=\{x:x \equiv 0 ~\mbox{on}~L\}$ with $|L |=n/k$,
then for any $x\in X$,
$$
\Pr(T(x)\in X) =\sum_t\Pr(|K\cap L |=t)2^{-(|L|-t)}
=\mu(X)\sum_t\Pr(|K\cap L |=t)2^t,
$$
and, since $|K\cap L|$ is essentially Poisson with mean 1,
the sum is
approximately
$ \eee^{-1}\sum_t2^t/t! = \eee$.
So Claim D fails for $\mu(X) = 2^{-n/k}$ and, as earlier,
it is %
natural to guess that it holds if $\mu(X)$ is much
%
greater  %
than this.

\mn
\emph{Proof of Claim D.}
Let $Q_r=\{y\in Q:|y|\leq r\}$.
The assumption on $\mu(X)$ implies that
$\mu(Q_{r-1})<\mu(X) \leq \mu(Q_r)$
for some $r >(1/2-o(1/k))n$, so Claim D follows from

\mn
\emph{Claim E.}
If $\vp =o(1/k)$ and $r > (1/2-\vp)n$, then for any $x\in Q$ and
$X\sub Q$ with
\beq{muX}
\mu(X) \leq \mu(Q_r),
\enq
we have

\beq{PrTx}
\Pr(T(x)\in X) < (1+o(1))\mu(Q_r).
\enq

\mn
\emph{Proof}.
We may assume $x=\underline{0}$, so that $\Pr(T(x)=y)$ is a decreasing
function of $|y|$.
We thus maximize $\Pr(T(x)\in X)$
subject to \eqref{muX} by taking
$
X=Q_r,
$
and \eqref{PrTx} is then a routine calculation
using
$$\mu(X) =\Pr({\rm Bin}(n,1/2)\leq r)$$
and
$$\Pr(T(x)\in X) =\Pr({\rm Bin}(n-k,1/2)\leq r)$$
(where ${\rm Bin}(\cdot,\cdot)$ denotes a binomially
distributed r.v.).\qed


\mn